\chapter{Relationship Label Definitions}\label{app:labels}

The definitions below were used both for ground-truth annotation and in the
prompt of the LLM classification method. $a$ denotes a provision from
\standardA{}, $b$ a provision from \standardB{}.

\section*{Equivalence}
$a$ and $b$ express the same security objective with comparable scope.
Compliance with one gives strong evidence of compliance with the other.
\emph{Decision aid:} if an auditor would accept evidence for $a$ as evidence
for $b$ without further argument, the pair is equivalent.

\section*{Subsumption (A broader)}
$a$'s objective fully contains $b$'s: everything $b$ requires is required by
$a$, but $a$ requires more. \emph{Decision aid:} compliance with $a$ implies
compliance with $b$, not conversely.

\section*{Subsumption (B broader)}
The mirror case: compliance with $b$ implies compliance with $a$.

\section*{Overlap}
The objectives intersect: part of what $a$ requires is also required by $b$,
but each provision also covers ground the other does not. Evidence can be
partially reused; a gap analysis is still needed for the uncovered parts.

\section*{Complementarity}
$a$ and $b$ address different aspects of a common security goal without
overlapping scope --- for example, one requires a technical mechanism and the
other requires user-facing guidance about it. Implementing both is what
produces the protection; neither substitutes for the other.

\section*{No relation}
No meaningful connection between the objectives. Used both for obviously
unrelated pairs and --- more importantly for evaluation --- for pairs that
share vocabulary but not objective.

\section*{Annotation rules}
Each pair receives exactly one label, a one-sentence justification, and a
confidence score: 1 (uncertain), 2 (reasonably confident), 3 (certain). When
several labels seem defensible, the annotator picks the most specific one
that the provision \emph{texts} (not assumed intent) support, and records the
doubt in the justification.
